\documentclass[a4paper, 14pt]{extarticle}

% Поля
%--------------------------------------
\usepackage{geometry}
\geometry{a4paper,tmargin=2cm,bmargin=2cm,lmargin=3cm,rmargin=1cm}
%--------------------------------------


%Russian-specific packages
%--------------------------------------
\usepackage[T2A]{fontenc}
\usepackage[utf8]{inputenc}
\usepackage[english, main=russian]{babel}

%--------------------------------------

\usepackage{textcomp}

% Красная строка
%--------------------------------------
\usepackage{indentfirst}               
%--------------------------------------             


%Graphics
%--------------------------------------
\usepackage{graphicx}
\graphicspath{ {./images/} }
\usepackage{float}%"Плавающие" картинки
%--------------------------------------

% Полуторный интервал
%--------------------------------------
\linespread{1.3}                    
%--------------------------------------

% Уменьшение расстояний между плавающими объектами
%--------------------------------------
\setlength{\floatsep}{10pt plus 2pt minus 2pt}
\setlength{\textfloatsep}{15pt plus 2pt minus 4pt}
\setlength{\intextsep}{10pt plus 2pt minus 2pt}
%--------------------------------------

%Выравнивание и переносы
%--------------------------------------
% Избавляемся от переполнений
\sloppy
% Запрещаем разрыв страницы после первой строки абзаца
\clubpenalty=10000
% Запрещаем разрыв страницы после последней строки абзаца
\widowpenalty=10000
%--------------------------------------

%Списки
\usepackage{enumitem}

\usepackage{minted}
%Подписи
\usepackage{caption} 

%Гиперссылки
% \usepackage{hyperref}

\usepackage[colorlinks, urlcolor = blue, filecolor = blue, citecolor = blue, linkcolor = black]{hyperref}


\hypersetup {
unicode=true
}

%Рисунки
%--------------------------------------
\DeclareCaptionLabelSeparator*{emdash}{~--- }
\captionsetup[figure]{labelsep=emdash,font=onehalfspacing,position=bottom}
%--------------------------------------

\usepackage{tempora}

%Листинги
%--------------------------------------
%Листинги
%--------------------------------------
\usepackage{listings}
\lstset{
  basicstyle=\ttfamily\footnotesize, 
  %basicstyle=\footnotesize\AnkaCoder,        % the size of the fonts that are used for the code
  breakatwhitespace=false,         % sets if automatic breaks shoulbd only happen at whitespace
  breaklines=true,                 % sets automatic line breaking
  captionpos=t,                    % sets the caption-position to bottom
  inputencoding=utf8,
  frame=single,                    % adds a frame around the code
  keepspaces=true,                 % keeps spaces in text, useful for keeping indentation of code (possibly needs columns=flexible)
  keywordstyle=\bf,       % keyword style
  numbers=left,                    % where to put the line-numbers; possible values are (none, left, right)
  numbersep=5pt,                   % how far the line-numbers are from the code
  xleftmargin=25pt,
  xrightmargin=25pt,
  showspaces=false,                % show spaces everywhere adding particular underscores; it overrides 'showstringspaces'
  showstringspaces=false,          % underline spaces within strings only
  showtabs=false,                  % show tabs within strings adding particular underscores
  stepnumber=1,                    % the step between two line-numbers. If it's 1, each line will be numbered
  tabsize=2,                       % sets default tabsize to 8 spaces
  title=\lstname                   % show the filename of files included with \lstinputlisting; also try caption instead of title
}
%--------------------------------------

%%% Математические пакеты %%%
%--------------------------------------
\usepackage{amsthm,amsfonts,amsmath,amssymb,amscd}  % Математические дополнения от AMS
\usepackage{mathtools}                              % Добавляет окружение multlined
\usepackage[perpage]{footmisc}


% графики
\usepackage{booktabs}        % Пакет для улучшенных таблиц
\usepackage{pgfplots}        % Пакет для построения графиков
\pgfplotsset{compat=1.17}    % Установка совместимости
\usepackage{tikz}            % Пакет для создания схем
\usetikzlibrary{positioning, arrows.meta}  % Библиотеки для TikZ

% Литература 

% переименовываем  список литературы в "список используемой литературы"
\addto\captionsrussian{\def\refname{Список литературы}}



\captionsetup[figure]{name=Рисунок, labelsep=space}


%--------------------------------------

%--------------------------------------
%			НАЧАЛО ДОКУМЕНТА
%--------------------------------------

\begin{document}

%--------------------------------------
%			ТИТУЛЬНЫЙ ЛИСТ
%--------------------------------------
\begin{titlepage}
\thispagestyle{empty}
\newpage


%Шапка титульного листа
%--------------------------------------
\vspace*{-60pt}
\hspace{-65pt}
\begin{minipage}{0.3\textwidth}
\hspace*{-20pt}\centering
\includegraphics[width=\textwidth]{emblem}
\end{minipage}
\begin{minipage}{0.67\textwidth}\small \textbf{
\vspace*{-0.7ex}
\hspace*{-6pt}\centerline{Министерство науки и высшего образования Российской Федерации}
\vspace*{-0.7ex}
\centerline{Федеральное государственное автономное образовательное учреждение }
\vspace*{-0.7ex}
\centerline{высшего образования}
\vspace*{-0.7ex}
\centerline{<<Московский государственный технический университет}
\vspace*{-0.7ex}
\centerline{имени Н.Э. Баумана}
\vspace*{-0.7ex}
\centerline{(национальный исследовательский университет)>>}
\vspace*{-0.7ex}
\centerline{(МГТУ им. Н.Э. Баумана)}}
\end{minipage}
%--------------------------------------

%Полосы
%--------------------------------------
\vspace{-25pt}
\hspace{-35pt}\rule{\textwidth}{2.3pt}

\vspace*{-20.3pt}
\hspace{-35pt}\rule{\textwidth}{0.4pt}
%--------------------------------------

\vspace{1.5ex}
\hspace{-35pt} \noindent \small ФАКУЛЬТЕТ\hspace{80pt} <<Информатика и системы управления>>

\vspace*{-16pt}
\hspace{47pt}\rule{0.83\textwidth}{0.4pt}

\vspace{0.5ex}
\hspace{-35pt} \noindent \small КАФЕДРА\hspace{50pt} <<Теоретическая информатика и компьютерные технологии>>

\vspace*{-16pt}
\hspace{30pt}\rule{0.866\textwidth}{0.4pt}

\vspace{11em}

\begin{center}
\Large {\bf Лабораторная работа №3} \\ 
\large <<Построение стохастическое модели в заданной предметной области>> \\
\large {по курсу <<Моделирование>>} \\

\end{center}\normalsize

\vspace{8em}

\vbox{%
\hfill%
\vbox{%
\hbox{Студент: Аринин М.А.}%
\hbox{Группа: ИУ9-82Б}%
\hbox{Преподаватель: Домрачева А.Б.}%
}%
} 

\bigskip

\vfill


\begin{center}
\textsl{Москва 2026}
\end{center}
\end{titlepage}
%--------------------------------------
%		КОНЕЦ ТИТУЛЬНОГО ЛИСТА
%--------------------------------------


\renewcommand{\ttdefault}{pcr}

\setlength{\tabcolsep}{3pt}

\newpage
\setcounter{page}{2}

\section{Постановка задачи}
Цель работы --- построить вычислительную модель, позволяющую по набору входных показателей относить объект к одному из двух классов. Наличие контрольного признака \texttt{Outcome} позволяет дополнительно проверить качество полученного разбиения, однако при построении модели этот признак не используется.

\textbf{Входные данные:}
\begin{itemize}
  \item файл \texttt{diabetes.csv} (\texttt{https://disk.yandex.ru/i/M35XVBrxn09\_iQ});
  \item числовые признаки объекта: \texttt{Pregnancies, Glucose, BloodPressure, SkinThickness, Insulin, BMI, DiabetesPedigreeFunction, Age}.
\end{itemize}

\textbf{Выходные данные:}
\begin{itemize}
  \item разбиение объектов на 2 класса (метка класса для каждого объекта);
  \item координаты центров классов в пространстве сниженной размерности;
  \item таблицы и графики для обоснования и интерпретации результатов.
\end{itemize}

\section{Построение математической модели}\label{sec:build_model}
\textbf{Тип модели.} Задача относится к методам многомерного статистического анализа. Рассматривается стохастическая модель, поскольку входные данные представлены выборкой случайных векторов (наборов показателей), а решение формируется на основе статистических оценок и эвристических критериев качества.

\subsection{Формализация задачи}
Пусть задана выборка объектов $X=\{x^{(1)},\dots,x^{(n)}\}$, где каждый объект $x^{(k)}$ описывается $p$ числовыми показателями. Требуется построить модель, которая сопоставляет каждому объекту метку класса $c^{(k)}\in\{1,2\}$ так, чтобы внутри класса объекты были близки по выбранной метрике, а между классами --- отличались.
\subsection{Бинаризация показателей}
Для каждого признака $x_i$ вводится событие
\[
A_i:\; x_i \le t_i,
\]
где $t_i$ --- порог бинаризации. В практической реализации используется порог по медиане $t_i=\mathrm{median}(x_i)$ как устойчивый к выбросам вариант, формирующий близкие по размеру подгруппы.

\subsection{Коэффициенты взаимосвязи}
После бинаризации для каждой пары признаков $(i,j)$ рассматриваются события $A_i:\;x_i\le t_i$ и $B_j:\;x_j\le t_j$.
Пусть для этой пары построена таблица сопряжённости 2$\times$2 со счётчиками
\[
\begin{aligned}
a&=n(A_i\cap B_j),\quad b=n(A_i\cap \overline{B_j}),\\
c&=n(\overline{A_i}\cap B_j),\quad d=n(\overline{A_i}\cap \overline{B_j}),\\
n&=a+b+c+d.
\end{aligned}
\]

В работе вычисляются несколько формальных коэффициентов взаимосвязи (без опоры на смысловую интерпретацию признаков):
\begin{itemize}
  \item коэффициент коллигации (lift)
  \[
  K_{ij}=\frac{P(A_i\mid B_j)}{P(A_i)},\qquad P(A_i)=\frac{a+b}{n},\qquad P(A_i\mid B_j)=\frac{a}{a+c};
  \]
  \item коэффициент контингенции Пирсона \(\varphi\) для 2$\times$2-таблицы
  \[
  \varphi_{ij}=\frac{ad-bc}{\sqrt{(a+b)(c+d)(a+c)(b+d)}};
  \]
  \item коэффициент ассоциации Юла \(Q\)
  \[
  Q_{ij}=\frac{ad-bc}{ad+bc};
  \]
  \item коэффициент коллигации Юла \(Y\)
  \[
  Y_{ij}=\frac{\sqrt{ad}-\sqrt{bc}}{\sqrt{ad}+\sqrt{bc}}.
  \]
\end{itemize}

Для дальнейшего сокращения размерности требуется «мера близости» признаков. Для симметричных коэффициентов (\(|\varphi|,|Q|,|Y|\)) она задаётся модулем. Для lift используется симметризация:
\[
S_{ij}=\max(K_{ij},K_{ji}).
\]

\subsection{Схема модели}
\begin{figure}[H]
  \centering
  \begin{tikzpicture}[
    node distance=1.6cm,
    box/.style={rectangle, draw, fill=#1!15, minimum width=3.8cm, minimum height=1cm, text centered, font=\footnotesize, text width=3.6cm},
    arrow/.style={->, >=stealth, thick}
  ]
    \node[box=blue] (data) {Исходные данные\\$\{x^{(k)}\}_{k=1}^n$};
    \node[box=green, below=of data] (bin) {Бинаризация\\$A_i: x_i \le t_i$};
    \node[box=yellow, below=of bin] (sim) {Оценка взаимосвязей\\$K_{ij},\ |\varphi_{ij}|,\ |Q_{ij}|,\ |Y_{ij}|$};
    \node[box=orange, below=of sim] (sel) {Выбор меры и порога\\перебор (grid search)\\по метрикам качества};
    \node[box=orange, below=of sel] (red) {Сокращение размерности\\группы зависимых\\признаков $\to$\\композитные $G_1..G_m$};
    \node[box=purple, below=of red] (km) {Кластеризация $k$-means\\($k=2$, евклидова метрика)};
    \node[box=gray, below=of km] (res) {Результат: метки классов\\и центры классов};

    \draw[arrow] (data) -- (bin);
    \draw[arrow] (bin) -- (sim);
    \draw[arrow] (sim) -- (sel);
    \draw[arrow] (sel) -- (red);
    \draw[arrow] (red) -- (km);
    \draw[arrow] (km) -- (res);
  \end{tikzpicture}
  \caption{Схема вычислительной модели классификации}
  \label{fig:model_scheme}
\end{figure}

\subsection{Сокращение размерности}
Сокращение размерности выполняется без жёсткого отбрасывания признаков. По выбранной мере взаимосвязи \(S_{ij}\) и порогу \(S_{\text{thr}}\) строится граф: вершины --- признаки, ребро \((i,j)\) проводится при \(S_{ij}\ge S_{\text{thr}}\). Далее выделяются компоненты связности (группы зависимых признаков), и каждая группа заменяется композитным показателем как среднее от стандартизованных признаков группы. После построения композитных признаков выполняется стандартизация:
\[
z=\frac{x-\mu}{\sigma}.
\]

\section{Решение вычислительных задач}
\subsection{Вычисление коэффициентов взаимосвязи}
Для каждой пары признаков $(i,j)$ после бинаризации строится таблица сопряжённости 2$\times$2 с элементами $(a,b,c,d)$. По ней вычисляются коэффициенты \(K_{ij}\), \(\varphi_{ij}\), \(Q_{ij}\), \(Y_{ij}\) (формулы приведены в разделе \ref{sec:build_model}). Ниже приведены визуализации матриц для используемых коэффициентов.

\begin{figure}[H]
  \centering
  \includegraphics[width=0.95\textwidth]{colligation_K_heatmap.png}
  \caption{Матрица lift-коэффициента коллигации \(K_{ij}=\frac{P(A_i\mid B_j)}{P(A_i)}\)}
  \label{fig:kheat}
\end{figure}

\begin{figure}[H]
  \centering
  \includegraphics[width=0.95\textwidth]{phi_heatmap.png}
  \caption{Матрица коэффициента контингенции Пирсона \(\varphi\) (\(r_B\)).}
  \label{fig:phiheat}
\end{figure}

\begin{figure}[H]
  \centering
  \includegraphics[width=0.95\textwidth]{yule_Q_heatmap.png}
  \caption{Матрица коэффициента Юла \(Q\)}
  \label{fig:yuleqheat}
\end{figure}

\begin{figure}[H]
  \centering
  \includegraphics[width=0.95\textwidth]{yule_Y_heatmap.png}
  \caption{Матрица коэффициента Юла \(Y\)}
  \label{fig:yuleyheat}
\end{figure}

\subsection{Порог и построение групп зависимых признаков}
По выбранным мере и порогу строятся группы зависимых признаков. Пример полученных групп приведён в табл.~\ref{tab:groups}.

\medskip
\textbf{Пояснение к агрегированному критерию выбора.}
Чтобы автоматически выбрать меру взаимосвязи и порог сокращения размерности (без подбора по \texttt{Outcome}), используется скаляризация многокритериальной задачи через эвристический score~\cite{podinovskiy2013}:
\[
\mathrm{score}=\mathrm{sil}-0.15\cdot\ln(1+\mathrm{DB})+0.02\cdot\ln(1+\mathrm{CH})-0.03\cdot m+0.10\cdot \mathrm{balance}.
\]
Здесь \(\mathrm{sil}\) (silhouette) и \(\mathrm{CH}\) (Calinski--Harabasz) должны расти, а \(\mathrm{DB}\) (Davies--Bouldin) --- уменьшаться~\cite{sklearn2026}; \(m\) штрафует избыточную сложность (число композитных признаков), \(\mathrm{balance}=\min(n_1,n_2)/\max(n_1,n_2)\) поощряет невырожденное разбиение по размерам кластеров. Логарифмирование \(\ln(1+x)\) применяется, чтобы привести метрики разного масштаба к сопоставимому влиянию на сумму (например, \(\mathrm{CH}\) может быть на порядки больше \(\mathrm{sil}\)). Весовые коэффициенты являются эвристическими и задают компромисс между качеством кластеризации, устойчивостью и простотой модели.

\medskip
\textbf{Замечание о мультиколлинеарности.}
В мультиколлинеарных пространствах (сильно коррелированные признаки) при построении композитных показателей меняется геометрия данных и масштабы расстояний; поэтому оптимизация внутренних метрик и простоты модели может давать предпочтение меньшему \(m\). Важно, что \textbf{score оптимизирует другое} — внутреннюю геометрию кластеров и простоту (через штраф за размерность), поэтому он может не совпасть с внешним критерием по \texttt{Outcome}.

\begin{table}[H]
\centering
\input{tables/reduction_groups.tex}
\caption{Группы признаков после сокращения размерности}
\label{tab:groups}
\end{table}

\subsection{Кластеризация методом $k$-means}
Кластеризация выполняется в пространстве композитных признаков $(G_1,\dots,G_m)$ по евклидовой метрике. Начальные центры выбираются как первая точка и наиболее удалённая от неё, затем итеративно выполняются шаги:
\begin{enumerate}
  \item назначить каждой точке ближайший центр;
  \item пересчитать центры как среднее по точкам кластера;
  \item повторять до стабилизации центров.
\end{enumerate}

Размеры кластеров:
\begin{table}[H]
\centering
\input{tables/cluster_sizes.tex}
\caption{Количество объектов в кластерах}
\end{table}

Центры кластеров (в композитных стандартизированных координатах):
\begin{table}[H]
\centering
\input{tables/cluster_centers.tex}
\caption{Центры кластеров в сниженной размерности}
\end{table}

\section{Практическая реализация}
\newenvironment{longlisting}{\captionsetup{type=listing}}{}
\begin{longlisting}
\caption{Листинг программной реализации (сокращённая версия для отчёта).}
\inputminted[frame=single,fontsize=\footnotesize,linenos,xleftmargin=1.5em,breaklines=true]{python}{modeling_process_lab3_report_clean.py}
\end{longlisting}

\section{Результаты}
Получено разбиение объектов на 2 класса. Ниже приведены результаты вычисления матриц взаимосвязей, выбора меры/порога для сокращения размерности и визуализация кластеризации.

\subsection{Выбор меры взаимосвязи и порога}
Мера взаимосвязи и порог для построения групп зависимых признаков выбираются автоматически перебором кандидатов: \(S_{ij}\in\{\max(K_{ij},K_{ji}),\ |\varphi_{ij}|,\ |Q_{ij}|,\ |Y_{ij}|\}\) и \(S_{\text{thr}}\) из заданной сетки. Каждый кандидат оценивается по внутренним метрикам качества кластеризации и балансу кластеров; при этом рассматриваются варианты размерности от \(m=2\) до \(m=p\) (где \(p\) --- число исходных признаков).

\textbf{Замечание о выборе.} Внутренние метрики (silhouette, Davies--Bouldin, Calinski--Harabasz) оценивают геометрию кластеров без учителя, поэтому «лучший» по ним вариант может не совпадать с разбиением, наиболее согласованным с контрольной меткой \texttt{Outcome}. Так как классы \texttt{Outcome} несбалансированы, метрика \(\mathrm{accuracy}\) может быть обманчивой; поэтому для внешней контрольной оценки используются \(\mathrm{balanced\ accuracy}\) и \(\mathrm{F1}\) (macro-F1). При расчёте внешних метрик выбирается лучшее сопоставление меток кластеров и \texttt{Outcome}.

\textbf{Выбранный вариант для итогового результата (по максимуму macro-F1):} \(S_{ij}=|Y_{ij}|\), \(S_{\text{thr}}=0.25\), итоговая размерность \(m=4\), \(\mathrm{F1_{macro}}=0.7261\), \(\mathrm{BalAcc}=0.7298\).

\textbf{Сравнение с вариантом, выбранным по score:} при \(S_{ij}=|Y_{ij}|\), \(S_{\text{thr}}=0.20\), \(m=2\) получаем \(\mathrm{F1_{macro}}=0.5842\), \(\mathrm{BalAcc}=0.5831\).

\begin{table}[H]
\centering
\scriptsize
\setlength{\tabcolsep}{2.5pt}
\resizebox{\textwidth}{!}{\input{tables/threshold_search.tex}}
\caption{Сравнение мер взаимосвязи и порогов по метрикам качества кластеризации и внешним контрольным метрикам (Acc, BalAcc, F1-macro)}
\end{table}

\begin{table}[H]
\centering
\input{tables/accuracy_by_m.tex}
\caption{Лучшие внешние метрики по \texttt{Outcome} в зависимости от числа композитных признаков \(m\)}
\end{table}

\begin{figure}[H]
  \centering
  \includegraphics[width=0.85\textwidth]{kmeans_best_score_2d.png}
  \caption{Разбиение $k$-means в 2D-проекции (вариант, выбранный по внутреннему score)}
  \label{fig:kmeans_best_score_2d}
\end{figure}

\begin{figure}[H]
  \centering
  \includegraphics[width=0.85\textwidth]{kmeans_best_macro_f1_2d.png}
  \caption{Разбиение $k$-means в 2D-проекции (вариант, выбранный по максимуму macro-F1 по \texttt{Outcome})}
  \label{fig:kmeans_best_macro_f1_2d}
\end{figure}

\begin{figure}[H]
  \centering
  \includegraphics[width=0.85\textwidth]{kmeans_best_macro_f1_3d.png}
  \caption{Разбиение $k$-means в 3D-проекции (вариант, выбранный по максимуму macro-F1; G1, G2, G3)}
  \label{fig:kmeans_best_macro_f1_3d}
\end{figure}

\section{Анализ моделей и проверка качества}
Качество кластеризации оценивается как внутренними метриками, так и внешней проверкой (без использования \texttt{Outcome} при построении модели).

\subsection{Внутренние метрики качества}
Для выбранного разбиения вычисляются стандартные внутренние метрики:
\begin{itemize}
  \item средний коэффициент силуэта (silhouette) --- чем больше, тем лучше;
  \item индекс Дэвиса--Болдина (Davies--Bouldin) --- чем меньше, тем лучше;
  \item индекс Калински--Харабаса (Calinski--Harabasz) --- чем больше, тем лучше.
\end{itemize}

\medskip
Согласно теории кластерного анализа, решение является принципиально неоднозначным: не существует единственного «наилучшего» критерия качества кластеризации, результат зависит от выбранной метрики расстояния и стандартизации признаков, а число кластеров часто задаётся исследователем. Поэтому в работе используются как внутренние метрики качества, так и внешняя контрольная оценка по \texttt{Outcome} (только для проверки результата, без участия в построении модели).

\subsection{Внешняя проверка по Outcome}
Столбец \texttt{Outcome} не использовался при построении модели и кластеризации. Однако его можно применить для внешней проверки согласованности разбиения: сравнивается распределение значений \texttt{Outcome} в полученных классах и вычисляются внешние контрольные метрики. Так как классы \texttt{Outcome} несбалансированы, для контроля используются \(\mathrm{BalAcc}\) и \(\mathrm{F1}\) (macro-F1).

\begin{table}[H]
\centering
\input{tables/accuracy.tex}
\caption{Внешние контрольные метрики по \texttt{Outcome} (лучшее сопоставление меток кластеров и \texttt{Outcome})}
\end{table}

\section{Выводы}
В ходе работы построена стохастическая вычислительная модель отнесения объектов к двум классам по входным данным из \texttt{diabetes.csv}. Получена полная цепочка вычислительных процедур:
\begin{itemize}
  \item выполнена бинаризация каждого числового признака по устойчивому порогу \(t_i=\mathrm{median}(x_i)\), что обеспечивает формирование сопоставимых по объёму подвыборок и снижает влияние выбросов;
  \item рассчитаны матрицы взаимосвязи признаков по нескольким формальным коэффициентам (\(K_{ij}\), \(\varphi\), Юла \(Q\), Юла \(Y\)), что позволяет не фиксировать заранее единственную меру зависимости;
  \item реализован алгоритм выбора меры взаимосвязи и порога сокращения размерности: перебираются кандидаты, строятся группы зависимых признаков и композитные показатели \(G_1,\dots,G_m\), после чего вариант оценивается по внутренним метрикам качества кластеризации и балансу кластеров для размерностей от \(m=2\) до \(m=p\) (где \(p\) --- число исходных признаков);
  \item выполнена кластеризация методом \(k\)-means в пространстве композитных признаков с получением центров кластеров и меток классов;
  \item проведена формальная интерпретация полученных классов через сравнение статистик исходных признаков внутри каждого класса.
\end{itemize}

Для контроля качества использована внешняя проверка по признаку \texttt{Outcome}. Из-за дисбаланса классов основная контрольная метрика --- macro-F1 (и balanced accuracy). При снятии ограничения на \(m\) оказалось, что вариант, оптимальный по внутреннему score, даёт меньшие внешние метрики (например, при \(m=2\) \(\mathrm{F1_{macro}}\approx 0.584\)), тогда как при \(m=4\) достигаются существенно лучшие значения (\(\mathrm{F1_{macro}}\approx 0.726\)). Поэтому в итоговых результатах выбран вариант, максимизирующий macro-F1, а внутренние метрики приведены как характеристика структуры кластеров.

\begin{thebibliography}{9}
\bibitem{podinovskiy2013}
Подиновский В.\,В., Потапов М.\,А. Метод взвешенной суммы критериев в анализе многокритериальных решений: \textit{pro et contra} // \textit{Бизнес-информатика}. 2013.

\bibitem{sklearn2026}
\textit{scikit-learn: руководство пользователя (русский перевод)}. Раздел «Кластеризация» и метрики качества (silhouette, Davies--Bouldin, Calinski--Harabasz). Дата обращения: 31.03.2026. 
URL: \url{https://scikit-learn.ru/stable/modules/clustering.html}.
\end{thebibliography}

\end{document}

